\documentclass[aspectratio=169]{beamer}

\usetheme{default}

% ===== Packages =====
\usepackage{amsmath}
\usepackage{amssymb}
\usepackage{booktabs}

% ===== 中文設定 =====
\usepackage{xeCJK}
\setCJKmainfont{PingFang TC}

% ===== Title =====
\title{YouBike Demand Around MRT Stations}
\subtitle{R for Data Science Final Project}
\author{Group 3}
\date{\today}

\begin{document}

\begin{frame}
\titlepage
\end{frame}

\begin{frame}{Variable Notation}

\small

\begin{columns}[T]

\column{0.48\textwidth}

\begin{itemize}
    \item $Y_{ist}$: YouBike trip count
    \item $Rain_{ist}$: Rain indicator
    \item $Temp_{ist}$: Temperature
    \item $Humidity_{ist}$: Relative humidity
\end{itemize}

\column{0.48\textwidth}

\begin{itemize}
    \item $Weekend_s$: Weekend indicator
    \item $Hour_t$: Hour fixed effects
    \item $Station_i$: MRT station fixed effects
    \item $\varepsilon_{ist}$: Error term
\end{itemize}

\end{columns}

\vspace{0.4cm}

\begin{center}
\footnotesize
$i$ = MRT station \qquad
$s$ = date \qquad
$t$ = hour
\end{center}

\end{frame}

%=
\begin{frame}{Model 0: Weather Only}

\[
Y_{ist}
=
\beta_0
+
\beta_1 Rain_{ist}
+
\beta_2 Temp_{ist}
+
\beta_3 Temp_{ist}^{2}
+
\beta_4 Humidity_{ist}
+
\varepsilon_{ist}
\]


\end{frame}
%
\begin{frame}{Model 1: Weather + Time}

\[
\begin{aligned}
Y_{ist}
=
\beta_0
+
\beta_1 Rain_{ist}
+
\beta_2 Temp_{ist}
+
\beta_3 Temp_{ist}^{2}
+
\beta_4 Humidity_{ist}
+
\beta_5 Weekend_s
+ \\
\sum_{h}
\gamma_h Hour_{ht}
+
\varepsilon_{ist}
\end{aligned}
\]

\begin{block}{Purpose}
Examine whether temporal patterns improve explanatory power.
\end{block}

\end{frame}
%
\begin{frame}{Model 2: Weekend $\times$ Hour}

\[
\begin{aligned}
Y_{ist}
=
\beta_0
+
\beta_1 Rain_{ist}
+
\beta_2 Temp_{ist}
+
\beta_3 Temp_{ist}^{2}
+
\beta_4 Humidity_{ist}
+
\beta_5 Weekend_s
+ \\
\sum_h \gamma_h Hour_{ht}
+
\sum_h \delta_h
(Weekend_s \times Hour_{ht})
+
\varepsilon_{ist}
\end{aligned}
\]

\begin{block}{Purpose}
Test whether weekday and weekend exhibit different intraday demand patterns.
\end{block}

\end{frame}
%
\begin{frame}{Model 3: Rain $\times$ Hour}

\[
\begin{aligned}
Y_{ist}
=
\beta_0
+
\beta_1 Rain_{ist}
+
\sum_h \delta_h
(Rain_{ist}\times Hour_{ht})
+
\beta_2 Temp_{ist}
+
\beta_3 Temp_{ist}^{2}
+
\beta_4 Humidity_{ist}
+ \\
\beta_5 Weekend_s
+
\varepsilon_{ist}
\end{aligned}
\]

\begin{block}{Purpose}
Test whether the impact of rain varies across different hours.
\end{block}

\end{frame}
%
\begin{frame}{Model 4: Station Fixed Effects}

\[
\begin{aligned}
Y_{ist}
=
\beta_0
+
\beta_1 Rain_{ist}
+
\beta_2 Temp_{ist}
+
\beta_3 Temp_{ist}^{2}
+
\beta_4 Humidity_{ist}
+
\beta_5 Weekend_s
+ \\
\sum_h \gamma_h Hour_{ht}
+
\sum_j \theta_j Station_{ij}
+
\varepsilon_{ist}
\end{aligned}
\]

\begin{block}{Purpose}
Control for station-specific demand differences.
\end{block}

\end{frame}
%
\begin{frame}{Model 5: Rain $\times$ Weekend}

\[
\begin{aligned}
Y_{ist}
=
\beta_0
+
\beta_1 Rain_{ist}
+
\beta_2 Weekend_s
+
\beta_3
(Rain_{ist}\times Weekend_s)
+
\beta_4 Temp_{ist}
+ 
\beta_5 Temp_{ist}^{2}
+ \\
\beta_6 Humidity_{ist}
+
\sum_h \gamma_h Hour_{ht}
+
\sum_j \theta_j Station_{ij}
+
\varepsilon_{ist}
\end{aligned}
\]

\begin{block}{Purpose}
Test whether the impact of rain differs between weekdays and weekends.
\end{block}

\end{frame}
%

%
\end{document}